\documentclass{standalone}
\usepackage{tikz}
\usepackage{tikz-3dplot}
\usepackage{graphicx}

\begin{document}
\tdplotsetmaincoords{-40}{0} % Set the viewpoint angles

\begin{tikzpicture}[tdplot_main_coords]
    
    \begin{scope}[canvas is xy plane at z=0]
        % \path[facestyle,shade] (0,0) rectangle (2,4);
        \coordinate (A) at (0, 0);
        \coordinate (B) at (4, 0);
        \coordinate (C) at (4, 3);
        \coordinate (D) at (0, 3);
        \draw[gray!50, fill=black!80] (A) -- (B) -- (C) -- (D) -- cycle;
    \end{scope}
    % Add the speckle effect (random small ellipses)
    \foreach \i in {1,...,800}
    {
        \pgfmathsetmacro{\cx}{rnd*3.85 + 0.075} % Random x-coordinate (range 0 to 3)
        \pgfmathsetmacro{\cy}{rnd*2.85 + 0.075} % Random y-coordinate (range 0 to 2)
        \pgfmathsetmacro{\rx}{rnd*0.05 + 0.01} % Random x-radius
        \pgfmathsetmacro{\ry}{rnd*0.05 + 0.01} % Random y-radius
        \pgfmathsetmacro{\rot}{rnd*360} % Random rotation angle (range 0 to 360)
        \pgfmathsetmacro{\gray}{rnd*0.9 + 0.1} % Random grayvalue

        
        % Draw ellipses only within the image plane boundaries
        \pgfmathsetmacro{\testx}{\cx + \rx}
        \pgfmathsetmacro{\testy}{\cy + \ry}
        
        \ifdim\testx pt<4mm
        \ifdim\testy pt<3mm
        % \tdplotsetrotatedcoords{0}{0}{0}
        % \tdplotsetrotatedcoordsorigin{\cx}{\cy}

        % \begin{scope}[tdplot_rotated_coords]
            \definecolor{ellipsecolor}{gray}{\gray}
            \begin{scope}[canvas is xy plane at z=0]
                \begin{scope}[xshift=28*\cx, yshift=28*\cy]
                    \draw[ellipsecolor, fill=ellipsecolor, rotate=\rot] (0,0) ellipse ({\rx} and {\ry});
                \end{scope}
            \end{scope}
        % \end{scope}
        \fi
        \fi
    }
    % \draw[fill=black] (rand,rand, 0) ellipse ({0.1} and {0.2});
    % \draw[fill=black] (rand,rand, 0) ellipse ({0.1} and {0.2});


    % Image plane coordinates
    
    % Camera coordinates
    \coordinate (camera) at (2, 1.5, 5);
    \coordinate (lens) at (2, 1.5, 3);
    
    % Draw the image plane (white or light gray)
    
    
    
    % Draw lines from camera to the image plane corners
    \draw[gray!50] (camera) -- (A);
    \draw[gray!50] (camera) -- (B);
    \draw[gray!50] (camera) -- (C);
    \draw[gray!50] (camera) -- (D);
    
    % Draw the lens of the camera
    % \draw[fill=gray!50] (lens) circle (0.5);
    
    % Draw the camera body
    \draw[white, fill=gray!70] (camera) -- (2.5, 1.5, 4.5) -- (1.5, 1.5, 4.5) -- cycle;
    \draw[white, fill=gray!70] (2.5, 1.5, 4.5) -- (2.5, 2, 4) -- (1.5, 2, 4) -- (1.5, 1.5, 4.5) -- cycle;
    \draw[white, fill=gray!70] (1.5, 1.5, 4.5) -- (1.5, 2, 4) -- (1.5, 2, 5) -- (1.5, 1.5, 5) -- cycle;
    \draw[white, fill=gray!70] (2.5, 1.5, 4.5) -- (2.5, 2, 4) -- (2.5, 2, 5) -- (2.5, 1.5, 5) -- cycle;
    \draw[white, fill=gray!70] (1.5, 1.5, 5) -- (2.5, 1.5, 5) -- (2.5, 2, 5) -- (1.5, 2, 5) -- cycle;
    
    % Add text labels
    % \node[above] at (A) {A};
    % \node[above] at (B) {B};
    % \node[right] at (C) {C};
    % \node[left] at (D) {D};
    % \node[above right, white, xshift=5mm] at (camera) {Camera};
    % \node[right] at (lens) {Lens};
    
    % Draw dashed line connecting camera and lens
    % \draw[dashed] (camera) -- (lens);
\end{tikzpicture}
\end{document}
